\subsection{Choice and Zorn's Lemma}

\begin{definition} \label{def:choice-function} \hlt{Choice Function}\\
Let $\mathcal{F}$ be a set of sets with $F\neq\varnothing$ for all $F\in\mathcal{F}$, and let $(X_i)_{i\in I}$ satisfy $X_i\neq\varnothing$ for all $i\in I$.
\begin{enumerate}[label=(\roman*)]
\setlength{\itemsep}{0pt}
\item A \hlt{choice function} on $\mathcal{F}$ is a map $c:\mathcal{F}\to\bigcup_{F\in\mathcal{F}}F$ such that $c(F)\in F$ for all $F\in\mathcal{F}$.
\item An \hlt{indexed choice function} for $(X_i)_{i\in I}$ is a map $c:I\to\bigcup_{i\in I}X_i$ such that $c(i)\in X_i$ for all $i\in I$.
\end{enumerate}
\end{definition}

\begin{principle} \label{princ:axiom-of-choice-analysis} \hlt{Choice Principle}\\
Let $(X_i)_{i\in I}$ be an indexed family with $X_i\neq\varnothing$ for all $i\in I$.\\
The principle asserts $\exists c:I\to\bigcup_{i\in I}X_i$ with $c(i)\in X_i$ for all $i\in I$; equivalently, $\prod_{i\in I}X_i\neq\varnothing$.\\
Construct $c$ when explicit; invoke choice only for arbitrary simultaneous choices.
\end{principle}

\begin{principle} \label{princ:zorns-lemma} \hlt{Zorn's Lemma}\\
Let $(X,\leq)$ be a partially ordered set with $X\neq\varnothing$. If every chain $C\subseteq X$ has an upper bound $u\in X$, meaning $c\leq u$ for all $c\in C$, then there exists $m\in X$ such that $m\leq x$ implies $x=m$ for all $x\in X$.
\end{principle}

\begin{definition} \label{def:inclusion-and-containment-orders-for-zorn} \hlt{Inclusion and Containment Orders}\\
Let $\mathcal{F}$ be a family of sets and let $\mathcal{A}\subseteq\mathcal{F}$.
\begin{enumerate}[label=(\roman*)]
\setlength{\itemsep}{0pt}
\item \hlt{Inclusion}: $A\leq B$ iff $A\subseteq B$.\\
An upper bound of $\mathcal{A}$ is $U\in\mathcal{F}$ with $A\subseteq U$ for all $A\in\mathcal{A}$\\
$M\in\mathcal{F}$ is maximal iff $M\subseteq B\in\mathcal{F}$ implies $B=M$.
\item \hlt{Reverse inclusion}: $A\leq B$ iff $B\subseteq A$.\\
An upper bound of $\mathcal{A}$ is $U\in\mathcal{F}$ with $U\subseteq A$ for all $A\in\mathcal{A}$;\\
$M\in\mathcal{F}$ is maximal iff $B\subseteq M$ with $B\in\mathcal{F}$ implies $B=M$.
\end{enumerate}
\end{definition}

\begin{method} \label{meth:zorn-maximality-template} \hlt{Zorn Maximality Template}\\
To prove existence by maximality, encode candidate objects as a poset $(P,\leq)$.
\begin{enumerate}[label=(\roman*)]
\setlength{\itemsep}{0pt}
\item Define $P$ so that a maximal $m\in P$ is a desired object. Then Prove $P\neq\varnothing$.
\item For an arbitrary chain $C\subseteq P$, construct $u\in P$ with $p\leq u$ for all $p\in C$.
\item Take a maximal $m\in P$ and verify the target property.
\end{enumerate}
For inclusion orders, the upper-bound candidate is often $\bigcup C$; for reverse-inclusion orders, it is often $\bigcap C$ or another object contained in every member of $C$. The key verification is always membership in $P$.
\end{method}
